\section{Introduction}
Speculative decoding saves sequential target invocations by verifying several draft tokens together~\cite{leviathan2023,chen2023}. That change also alters the target's execution geometry. Floating-point evaluation can depend on query width, so mathematically equivalent batching can change the greedy response of a deployed quantized model. When an application must retain its existing regression baseline, recovering the reference output becomes part of the optimization problem.

The decisive choice is where to impose that numerical constraint. Global serialization (GS) evaluates one target row at a time through projections, attention and feed-forward layers. It recovers a serial execution organization at the cost of surrounding batching. Per-query attention (PQ) confines the separation to attention but constructs a subgraph for each query. We ask whether the same local organization can be expressed more cheaply, and whether the resulting repair changes the economic case for drafting.

Phase-scoped query-to-batch attention (QB) answers the first question through graph lowering. It views Q as $(d,1,h,q)$ instead of $(d,q,h,1)$ and transforms the mask in parallel. Existing attention operators process the queries as separate batch instances while K/V storage remains shared within one sequence. Projections and feed-forward computation stay batched. Role/phase guards and a layout-sensitive graph cache bind this representation to the native speculative transaction, including execution with observation disabled.

The contribution is this runtime contract and its measured consequences. Query independence itself is standard attention structure. Batch-invariant kernels control reductions~\cite{thinkingmachines}; CoRun controls prefill and decode shapes~\cite{corun}; LLM-42 verifies a fast path under a fixed-shape schedule~\cite{llm42}. QB instead preserves a local one-query attention organization through views, head mapping, phase dispatch and graph reuse, without adding a numerical-consistency verifier. Section~\ref{sec:related} compares these control points explicitly.

The matched CUDA experiments separate three decisions. First, local repair changes drafting from a net cost to a net benefit for Qwen3 and Llama; both PQ and QB cross this boundary. Second, QB lowers the representation cost within local repair: six sessions across two dates yield speculative PQ/QB request-time ratios of 1.032, 1.066 and 1.029 for Qwen3, SmolLM2 and Llama. A 72-context replay finds bitwise agreement at QK, softmax and AV. Third, a cheaper representation does not guarantee profitable drafting: SmolLM2 benefits from QB but remains faster without the draft.

Transfer and deployment test the same argument at larger boundaries. Vulkan state interventions show that endpoint-equivalent states can contain noninterchangeable KV components. A grouped-query port preserves the existing AV evaluation while dispatching independent queries. Finally, mapping and fresh-context service studies measure the complete transaction~\cite{mappingrecords,b570servicerecords,rtxservicerecords}. Neither fresh-context panel establishes a practical QB/PQ increment. Their different net draft costs reinforce the need to compare qualified policies at the intended request boundary.

The contributions are:
\begin{enumerate}[label=(\arabic*),itemsep=2pt]
\item A phase-scoped attention transformation with paired Q/mask views, shared-KV and head semantics, ordered output recovery and layout-sensitive graph reuse.
\item Matched evidence that repair scope changes draft profitability, followed by a separate measurement of QB's incremental cost benefit over PQ.
\item State and backend transfer tests, together with measurements of the representation benefit at the deployment request boundary.
\end{enumerate}
The supplement supplies the full comparisons, failures and implementation records; ancillary files pin their source identities~\cite{studyrecords}.
